% Manuale Utente - Progetto Ruggine
\documentclass[a4paper,12pt]{book}
\usepackage[utf8]{inputenc}
\usepackage[italian]{babel}
\usepackage{hyperref}
\usepackage{graphicx}

\title{Manuale Utente}
\author{Gruppo 50}
\date{\today}

\begin{document}
% Frontespizio
\begin{titlepage}
    \centering
    \includegraphics[width=0.3\textwidth]{img/ruggine_logo.png}\\
    {\Huge\bfseries Manuale Utente\par}
    \vspace{0.5cm}
    {\Large Progetto \textbf{Ruggine}\par}
    \vspace{0.5cm}
    {\large Corso di Programmazione di Sistema \\ 2025/2026\par}
    \vspace{1.5cm}
    {\large\textbf{Gruppo 50}}\\[0.5cm]
    Alessandro Annecchini \\ Riccardo Marconi \\ Francesco Palmisani \\ Giosuè Pinto
    \vfill
    {\large \today\par}
\end{titlepage}

\frontmatter
\tableofcontents

\mainmatter
\chapter{Introduzione}
Questo manuale guida l'utente nell'utilizzo della web app Ruggine, illustrando le principali funzionalità disponibili.

\chapter{Panoramica dell'applicazione}
Ruggine è una piattaforma di chat che permette di comunicare in tempo reale, creare gruppi, inviare e ricevere inviti.

\chapter{Accesso e Registrazione}
\section{Come registrarsi}
\section{Come effettuare il login}
\section{Recupero password (se previsto)}

\chapter{Chat e Messaggistica}
\section{Inviare e ricevere messaggi}
\section{Chat di gruppo e chat dirette (DM)}
\section{Notifiche e aggiornamenti in tempo reale}

\chapter{Gestione Gruppi}
\section{Creare un nuovo gruppo}
\section{Aggiungere/rimuovere membri}
\section{Visualizzare la lista dei gruppi}

\chapter{Inviti}
\section{Inviare un invito a un altro utente}
\section{Accettare o rifiutare un invito}
\section{Visualizzare inviti ricevuti e inviati}

\chapter{FAQ e Risoluzione Problemi}
\section{Problemi di accesso}
\section{Messaggi non inviati}
\section{Chi contattare per assistenza}

\end{document}
